\documentclass[aspectratio=169]{beamer}

\usetheme{default}
\usecolortheme{seahorse}
\setbeamertemplate{navigation symbols}{}
\setbeamertemplate{footline}[frame number]

\usepackage{graphicx}
\usepackage{booktabs}

\graphicspath{{figures/}}

\title{MAC Detector Geometry: Four Cases}
\subtitle{Comparing detector mounting strategies for HRD}
\date{\today}

\begin{document}

\begin{frame}
  \titlepage
\end{frame}

% ---------------------------------------------------------------------------
\begin{frame}{Cases considered}
  \begin{description}
    \item[Case A] \textbf{N independent rotating arms.} Each detector arm
      tracks its crystal's reflected beam.  Crystal-to-detector distance and
      incidence angle are constant by construction.  Requires one rotation
      motor per detector.

    \item[Case B] \textbf{Rigid annular plate rotating about the sample.}
      Detectors sit on a circle of fixed radius $D = R + R_d$ centred on the
      sample.  A single rotation motor suffices; both $r_d$ and incidence
      angle vary with energy.

    \item[Case C] \textbf{Flat plate rotating about the central crystal.}
      The plate normal tracks the central crystal's diffracted beam.
      Outer crystals hit the plate off-normal by their arm offset (constant).
      $r_d$ varies with energy.

    \item[Case D] \textbf{Fixed flat plate rotating about the sample.}
      The plate is fixed perpendicular to the central arm direction at
      $R + R_d$.  Both $r_d$ and incidence angle vary with energy for all
      crystals.
  \end{description}
\end{frame}

% ---------------------------------------------------------------------------
\begin{frame}{How to read the diagnostic panels}
  Each case figure has three columns:
  \begin{columns}
    \begin{column}{0.33\textwidth}
      \textbf{Real-space layout}\\[4pt]
      Crystal positions (coloured dots, magma by index).
      Detector positions (scatter, viridis by energy).
      Reference circles and plate guide lines where applicable.
    \end{column}
    \begin{column}{0.33\textwidth}
      \textbf{Crystal–detector metrics}\\[4pt]
      \emph{Top:} crystal-to-detector distance $r_d$ vs energy.\\
      \emph{Bottom:} beam incidence angle on detector face vs energy.\\
      Flat lines indicate Case~A's ideal behaviour.
    \end{column}
    \begin{column}{0.33\textwidth}
      \textbf{Rigid-body analysis}\\[4pt]
      \emph{Top:} per-detector residual after best-fit rigid transform.
      Near-zero $\Rightarrow$ one motor suffices.\\
      \emph{Middle:} optimal rotation angle vs energy.\\
      \emph{Bottom:} optimal translation magnitude vs energy.
    \end{column}
  \end{columns}

  \vspace{6pt}
  \textbf{Rigid-body transform} (Kabsch/SVD): at each energy find the
  rotation $R$ and translation $t$ that best maps the lowest-energy detector
  configuration onto the current one.  Residuals are what is left over.
\end{frame}

% ---------------------------------------------------------------------------
\begin{frame}{Case A — independent rotating arms}
  \begin{center}
    \includegraphics[width=\linewidth]{mac_case_A}
  \end{center}
\end{frame}

% ---------------------------------------------------------------------------
\begin{frame}{Case A — notes}
  \begin{itemize}
    \item $r_d = R_d$ exactly, incidence angle $= 0°$ for all crystals and
      energies — the detector always faces the beam.
    \item Rigid-body residual $< 1$\,mm: the detector array motion is nearly
      (but not exactly) a rigid rotation because each arm pivots about a
      different crystal position.
    \item Rotation angle $\theta \lesssim 1.4°$ over the full energy range;
      translation $< 5$\,mm.
    \item \textbf{Cost:} one rotation motor per arm ($N = 12$ motors total).
  \end{itemize}
\end{frame}

% ---------------------------------------------------------------------------
\begin{frame}{Case B — rigid annular plate}
  \begin{center}
    \includegraphics[width=\linewidth]{mac_case_B}
  \end{center}
\end{frame}

% ---------------------------------------------------------------------------
\begin{frame}{Case B — notes}
  \begin{itemize}
    \item Detectors on a circle of radius $D = R + R_d$ about the sample.
    \item Rigid-body residual $= 0$ exactly: the motion \emph{is} a pure
      rotation about the sample origin.  Translation $= 0$, $\theta$ up to
      $1.4°$.
    \item $r_d$ varies $\sim 5$\,mm and incidence angle varies
      $\sim 12°$ across the energy range.
    \item \textbf{Cost:} one rotation motor for the whole plate.
      The varying incidence angle may reduce detector efficiency at the
      extremes of the energy range.
  \end{itemize}
\end{frame}

% ---------------------------------------------------------------------------
\begin{frame}{Case C — flat plate rotating about the central crystal}
  \begin{center}
    \includegraphics[width=\linewidth]{mac_case_C}
  \end{center}
\end{frame}

% ---------------------------------------------------------------------------
\begin{frame}{Case C — notes}
  \begin{itemize}
    \item Plate normal tracks the central crystal's diffracted beam.
      Incidence angle $= $ arm offset (constant per crystal, $\pm 11°$
      for outermost).
    \item Rigid-body residual up to $\sim 12$\,mm: the flat plate geometry
      cannot be captured by a single rigid motion.
    \item Optimal rotation $\theta$ up to $13.3°$ (= twice the Bragg angle
      sweep); translation up to $220$\,mm because the array centroid orbits
      \texttt{rot\_center\_C}.
    \item \textbf{Cost:} one rotation motor for the plate, but independent
      in-plane motion per detector is required to keep the beam on each
      detector.
  \end{itemize}
\end{frame}

% ---------------------------------------------------------------------------
\begin{frame}{Case D — fixed flat plate about the sample}
  \begin{center}
    \includegraphics[width=\linewidth]{mac_case_D}
  \end{center}
\end{frame}

% ---------------------------------------------------------------------------
\begin{frame}{Case D — notes}
  \begin{itemize}
    \item Plate fixed perpendicular to the central arm direction at
      distance $R + R_d$ from the sample.
    \item No rotation of the plate with energy ($\theta \approx 0$);
      translation up to $\sim 29$\,mm as the beam sweeps across the plate.
    \item Rigid-body residual up to $\sim 6$\,mm: smaller than Case~C but
      still requires independent per-detector motion.
    \item Incidence angle varies with energy for all crystals (including
      the central one), unlike Case~C.
    \item \textbf{Cost:} no rotation motor needed, but in-plane motion
      per detector is still required.
  \end{itemize}
\end{frame}

% ---------------------------------------------------------------------------
\begin{frame}{Summary comparison}
  \centering
  \begin{tabular}{lcccc}
    \toprule
    & \textbf{A} & \textbf{B} & \textbf{C} & \textbf{D} \\
    \midrule
    $r_d$ constant?              & \checkmark & $\times$ & $\times$ & $\times$ \\
    Incidence angle constant?    & \checkmark & $\times$ & per crystal & $\times$ \\
    Rigid-body residual (mm)     & $<1$       & $0$      & $\sim 12$ & $\sim 6$ \\
    Motors (plate-level)         & 0          & 1 rot    & 1 rot     & 0        \\
    Motors (per detector)        & 1 rot      & 0        & 1 in-plane & 1 in-plane \\
    \bottomrule
  \end{tabular}

  \vspace{8pt}
  \small Case~B is the unique case where a single rigid-body motion exactly
  captures all detector positions at all energies.  Cases~C and~D require
  per-detector in-plane motion; Case~A requires per-detector rotation.
\end{frame}

\end{document}
